\scanpage{12}
\mspage{26}

\uline{Théorème} \inserted{4.}6. Soient $A$ un anneau préadmissible muni d'un idéal de
définition $I=I_0$ et d'un système fondamental $(I_i)$ d'idéaux de définition
$I_0$, $B$ une $A$-algèbre topologique. On suppose\deleted{\unsure{x}} les $A_i=A/I_iA$ noethériens,
et de plus les conditions suivantes sur $B$:

\deleted{\unsure{a) Les} $\unsure{A_i=A/I_i}$ \unsure{noethériens}}

\unsure{a)} Les module\deleted{s} topologique\deleted{s}
\deleted{$\unsure{B_i=B/I_iB}$} \inserted{$B$ sur $A$}
\deleted{\unsure{sont}} satisf\deleted{ont}\inserted{ait} la condition P.

\inserted{b) $B_0$ est formellement simple sur $A_0$}

\deleted{b)}\inserted{c)} $B$ est adique \deleted{pour un idéal contenant $IB$}

\inserted{d)} Les idéaux $I_iB$ dans $B$ sont fermés (\deleted{inutile}\inserted{automatiquement vérifié} si $B$ est noethérien).

Sous ces conditions, on a ce qui suit\deleted{\unsure{x}}:

(i) $B$ est formellement simple sur $A$.

(ii) Soit $C$ une $A$-algèbre topologique satisfaisant à \unsure{c)} et d)
(mais pas nécessairement à a) ), et soit $u_0:B_0\xrightarrow{\sim}C_0$ un isomorphisme
d'algèbres topologiques. Alors il existe un $A$-homomorphisme continu
$u:B\longrightarrow C$ se réduisant suivant $u_0$ (en vertu de (i) et de 1.7.). Ce
dernier étant choisi, les conditions suivantes sont équivalentes: \deleted{\unsure{1°}}

\deleted{\unsure{a)}}\inserted{1)} $B$ est formellement simple sur $A$.

2) \deleted{$B$ est} $B$ satisfait la condition a) ci-dessus, i.e. \deleted{est} satisfait
la condition P.

3) Les $B_i$ sont plats sur $A_i$.

4) L'homomorphisme $u:B\longrightarrow C$ est un isomorphisme\deleted{\unsure{,}} de $A$-algèbres top.

\inserted{5) $\unsure{C}$ est \unsure{isom.} à $B$ comme $A$-algèbre top.}

Démonstration. L'assertion (i) résulte déjà de a)\inserted{et b)}, grâce à 3.1.
Prouvons donc (ii). En vertu de (i) déjà prouvé, on a $2)\Longrightarrow1)$, et
évidemment $4)\Longrightarrow2)$, d'ailleurs $2)\Longrightarrow3)$ en vertu de 3.5., grâce au
fait que les $A_i$ sont noethériens. Reste à prouver $1)\Longrightarrow4)$ et $3)\Longrightarrow4)$.

Supposons 1) vérifié, alors en vertu de 1.7. $u_0^{-1}:C_0\longrightarrow B_0$
se prolonge en un homomorphisme d'algèbres top $v:C\longrightarrow B$ (grâce aux
hypothèses \unsure{c)} et d) ). On a déjà noté que les $B_i$ sont plats sur les
$A_i$. Le lemme \inserted{4.}2. (i) et (iii) (où on échange les rôles de $B,C$)

\mspage{11}
